\chapter{Le mobile : Android, iOS, HarmonyOS}

Le mobile ajoute quelque chose que le bureau n'a pas : un \textbf{manifeste}. Il
ne suffit plus de compiler et de lier --- il faut déclarer une identité, des
permissions, une orientation, et signer. Ce chapitre couvre les trois systèmes.

\section{Android}

\begin{nkcode}{Applications/NkDames/NkDames.jenga -- reel, abrege}
with filter("system:Android"):
    windowedapp()
    usetoolchain("android-ndk")
    androidapplicationid("com.nkentseu.nkdames")
    androidisgame(True)
    androidminsdk(24)
    androidtargetsdk(34)
    androidcompilesdk(34)
    androidabis(["armeabi-v7a", "arm64-v8a", "x86", "x86_64"])
    androidnativeactivity(True)
    androidscreenorientation("portrait")
    androidallowrotation(False)
    androidstl("c++_shared")
    links(["android", "log", "EGL", "GLESv3", "dl", "OpenSLES"])
\end{nkcode}

\begin{center}
\begin{tabular}{@{}ll@{}}
\toprule
\textbf{Fonction} & \textbf{Ce qu'elle décide} \\
\midrule
\texttt{androidapplicationid} & l'identité du paquet --- unique sur l'appareil \\
\texttt{androidminsdk} & la version d'Android la plus ancienne acceptée \\
\texttt{androidabis} & les architectures embarquées \\
\texttt{androidnativeactivity} & pas de code Java du tout \\
\texttt{androidscreenorientation} & ce que le manifeste impose \\
\texttt{androidstl} & quelle bibliothèque standard C++ embarquer \\
\bottomrule
\end{tabular}
\end{center}

\begin{nkretenir}
\textbf{\texttt{androidabis} multiplie la taille du paquet.} Quatre
architectures, c'est quatre bibliothèques natives dans l'APK. Le paquet
« universel » d'un jeu de ce dépôt pèse 7 Mo pour 2 Mo sur le bureau.

On les garde pour un APK de test qui s'installe partout ; pour une distribution,
on découpe.
\end{nkretenir}

\begin{nkretenir}[Une NativeActivity n'a pas de \texttt{classes.dex}, et c'est normal]
\texttt{androidnativeactivity(True)} produit un manifeste avec
\texttt{android:hasCode="false"} : il n'y a \textbf{aucun} code Java, donc rien à
convertir en \texttt{dex}.

\alerte{} Un contrôle qui exige un \texttt{classes.dex} déclarera donc cassé un APK qui
s'installe et fonctionne. \emph{Vérifié : les APK de ce dépôt n'en ont pas, et
tournent sur un téléphone réel.}
\end{nkretenir}

\begin{nkpiege}[Sans orientation déclarée, l'image sort tournée]
Le système impose la sienne. Le descripteur commente d'ailleurs son choix :
« PORTRAIT : c'est l'orientation du jeu [\dots] ; la mise en page SAIT faire les
deux ; c'est un choix de produit, pas une limite technique ».

Et l'on vérifie ce que le paquet porte \emph{réellement}, pas ce que le
descripteur demande :

\begin{nkcode}{aapt dump xmltree -- capture reelle}
A: android:screenOrientation(0x0101001e)=(type 0x10)0x1
\end{nkcode}

\texttt{0x1} vaut \texttt{SCREEN\_ORIENTATION\_PORTRAIT}.

\alerte{} \textbf{Et cherchez la valeur, pas le mot.} Dans un manifeste binaire,
\texttt{screenOrientation} est un \textbf{entier} : chercher la chaîne
« portrait » rend ABSENT sur un paquet parfaitement verrouillé. Mesuré --- et
j'ai failli annoncer un défaut sur cet instrument aveugle.
\end{nkpiege}

\begin{nkretenir}[Ce qu'il faut sur la machine]
\begin{center}
\begin{tabular}{@{}ll@{}}
\toprule
\texttt{ANDROID\_SDK\_ROOT} & le SDK \\
\texttt{ANDROID\_NDK\_ROOT} & le NDK \\
\texttt{JAVA\_HOME} & un JDK 17 --- pour empaqueter et signer \\
\texttt{\textasciitilde/.android/debug.keystore} & la clé de test \\
\bottomrule
\end{tabular}
\end{center}

\textbf{Les deux dernières lignes sont celles qu'on oublie}, et leur absence ne
casse pas la compilation : elle produit un APK \emph{sans signature} que
\texttt{adb install} refuse en silence.

Vérification rapide d'un APK : il doit contenir \texttt{META-INF/MANIFEST.MF} et
un \texttt{META-INF/*.RSA}.
\end{nkretenir}

\section{iOS}

\begin{nkcode}{Applications/NkDames/NkDames.jenga -- reel}
with filter("system:iOS"):
    windowedapp()
    frameworks(["UIKit", "QuartzCore", "OpenGLES", "AVFoundation", "Metal",
                "AudioToolbox", "CoreAudio"])
    iosbundleid("com.nkentseu.nkdames")
    iosversion("1.0")
    iosminsdk("14.0")
    iosbuildnumber(1)
\end{nkcode}

\begin{nkpiege}[La liste des frameworks n'est PAS celle de macOS]
\textbf{\texttt{AudioUnit} n'existe pas sur iOS.} C'est un framework distinct sur
macOS seulement ; sur iOS, ses symboles --- \texttt{AudioComponent*},
\texttt{AudioOutputUnit*} --- vivent dans \texttt{AudioToolbox}.

Recopier la liste macOS donne \texttt{ld: framework AudioUnit not found}. Mesuré
sur l'intégration continue iOS.

\emph{Un précédent se vérifie sur SA plateforme.} C'est vrai des frameworks
Apple, et c'est vrai de bien d'autres choses.
\end{nkpiege}

\begin{nkretenir}[Ce que le descripteur avoue lui-même]
% `escapeinside` : le panneau ne s'imprime pas dans la police du code, et la
% regle `literate` du preambule ne mord pas dans un bloc de code -- mesure faite.
\begin{nkcode}[listing options={escapeinside={(*}{*)}}]{Applications/NkDames/NkDames.jenga}
# (*\alerte{}*) Non verifie sur machine : la construction iOS exige macOS + Xcode.
# Ces lignes suivent le modele de NKPA, qui est dans le meme cas. Elles
# rendent la cible DECLAREE, pas EPROUVEE.
\end{nkcode}

\textbf{C'est la distinction que ce livre tient partout.} Une cible déclarée a
des lignes dans un descripteur ; une cible éprouvée a produit un artefact qu'on a
lancé. Les confondre fait annoncer un support qu'on n'a pas.

Le même commentaire note que le filtre iOS \textbf{n'existait pas} jusqu'au
2026-08-27, alors que l'en-tête du fichier annonçait « Mobile : Android, iOS,
HarmonyOS » --- \emph{une plateforme qu'on croit couverte parce qu'un commentaire
la nomme}.
\end{nkretenir}

\section{HarmonyOS}

\begin{nkcode}{Applications/NkDames/NkDames.jenga -- reel}
with filter("system:HarmonyOS"):
    windowedapp()
    usetoolchain("ohos-ndk")
    harmonybundlename("com.nkentseu.nkdames")
    harmonyminsdk(12)
    harmonytargetapi(12)
    harmonyorientation("portrait")
    links(["hilog_ndk.z", "EGL", "GLESv3", "vulkan"])
\end{nkcode}

\begin{nkretenir}
La structure est celle d'Android, les noms changent :
\texttt{harmonybundlename} pour l'identité, \texttt{hilog\_ndk.z} pour le
journal --- l'équivalent de \texttt{log}.

\textbf{Et l'orientation s'y déclare pour la même raison} : sans elle, le
système impose la sienne et l'image sort tournée d'un quart de tour.
\end{nkretenir}

\begin{nkpiege}
\textbf{\texttt{jenga package --platform harmonyos} exige un
\texttt{hvigor-config.json5}} que les applications de ce dépôt n'ont pas. Leur
\texttt{.hap} vient donc de \texttt{jenga build}, pas de \texttt{package}.

C'est écrit dans le script d'empaquetage, et c'est le genre de chose qu'il vaut
mieux découvrir dans un livre que dans un journal d'erreurs.
\end{nkpiege}

\section{Le point commun des trois}

\begin{nkretenir}
\textbf{Sur mobile, la construction n'est que la moitié du travail.} Il reste à
produire un manifeste juste, à embarquer les bonnes architectures, et à signer.

Et ces trois-là échouent différemment de la compilation : \textbf{ils ne
produisent pas d'erreur de compilateur.} L'APK se construit, il ne s'installe
pas ; le manifeste se génère, il déclare la mauvaise orientation ; le paquet
existe, il n'est pas signé.

\emph{La vérification d'un artefact mobile porte donc sur le PAQUET, jamais sur
la sortie de la construction.}
\end{nkretenir}

\begin{nkbilan}
Le mobile ajoute un manifeste : identité, permissions, orientation, signature.
Android se décrit par une quinzaine d'appels dont l'oubli le plus coûteux est
l'orientation ; ses APK NativeActivity n'ont légitimement pas de
\texttt{classes.dex}, et il faut vérifier la valeur dans le manifeste binaire, pas
le mot. iOS partage la forme de macOS mais pas sa liste de frameworks ---
\texttt{AudioUnit} n'y existe pas --- et le descripteur avoue être déclaré sans
être éprouvé. HarmonyOS reprend la structure d'Android sous d'autres noms.
Retenez surtout que ces trois systèmes échouent \emph{après} la compilation, donc
que la vérification porte sur le paquet.
\end{nkbilan}

\begin{nkquestions}
\begin{enumerate}
\item Que produit \texttt{androidnativeactivity(True)} dans le manifeste, et
      quelle conséquence sur le contenu de l'APK ?
\item Comment vérifier l'orientation réellement déclarée par un APK ?
\item Pourquoi \texttt{AudioUnit} casse-t-il la construction iOS ?
\item Quelle différence entre une cible \emph{déclarée} et une cible
      \emph{éprouvée} ?
\item Qu'est-ce qui manque à un APK quand \texttt{JAVA\_HOME} n'est pas défini ?
\end{enumerate}
\end{nkquestions}

\begin{nkpratique}
Prenez un APK --- le vôtre, ou l'un de ceux du dépôt --- et établissez quatre
faits par la commande, pas par le descripteur :

\begin{enumerate}
\item son identité de paquet ;
\item son orientation déclarée ;
\item les architectures qu'il embarque ;
\item s'il est signé.
\end{enumerate}

\texttt{aapt dump badging} et \texttt{aapt dump xmltree} répondent aux trois
premières ; la quatrième se lit dans la liste des fichiers du zip.
\end{nkpratique}

\begin{nkdemo}
Prouvez qu'un manifeste binaire ne se lit pas comme du texte.

\begin{enumerate}
\item Cherchez la chaîne \texttt{portrait} dans \texttt{AndroidManifest.xml}
      extrait d'un APK verrouillé en portrait. Notez le résultat.
\item Faites lire le même fichier par \texttt{aapt dump xmltree}. Notez la valeur
      de \texttt{screenOrientation}.
\end{enumerate}

\textbf{Ce que la démonstration établit} : la première méthode dit ABSENT sur un
paquet parfaitement correct. Une valeur d'énumération y est un \emph{entier}.

\alerte{} Et posez le contrôle positif : cherchez une chaîne dont vous savez qu'elle est
là --- \texttt{NativeActivity}, par exemple. Si elle non plus ne sort pas, votre
lecteur est aveugle et le premier résultat ne prouvait rien.
\end{nkdemo}

\begin{nklogique}
Votre APK se construit, mais \texttt{adb install} le refuse sans message clair.

\begin{enumerate}
\item Énumérez quatre causes possibles, en distinguant celles qui viennent de la
      construction de celles qui viennent de l'appareil.
\item Pour chacune, dites ce que vous regarderiez \emph{dans le paquet} --- et
      non dans la sortie de la construction.
\item Une cause ne se voit ni dans le paquet ni dans la construction. Laquelle,
      et où la prendriez-vous ?
\end{enumerate}
\end{nklogique}
